\providecommand{\EqOverGt}{\mathrel{\substack{=\\>}}}
\SourceLeaf{492}

\begin{center}
\resizebox{.60\linewidth}{!}{\input{diagrams/native492_pentagon.tex}}
\end{center}
\[
\begin{gathered}\cos A=\alpha,\quad \cos B=\beta,\quad \cos C=\gamma,\\[9pt]\quad{}\quad \cos D=\delta,\quad \cos E=\varepsilon\end{gathered}
\]

Um die Eckpunkte in ganzen complexen Zahlen ausgedrückt zu erhalten,
seien $p,p',p'',p''',p''''$ fünf complexe ganze Zahlen, und zwar
\[
p=a+bi,\quad p'=a'+b'i,\quad\text{etc.}
\]
Man setze
\[
a'a'''+b'b'''=(1,3),\quad a''a''''+b''b''''=(2,4)\quad\text{u.\,s.\,w.}
\]
und nehme für die Eckpunkte
\par\nopagebreak% portrait-table start
\begingroup\small\setlength{\arraycolsep}{2pt}\[
\begin{array}{rccccc} & (1,3)(2,4)p, & (2,4)(3,0)p', & (3,0)(4,1)p'', & (4,1)(0,2)p''', & (0,2)(1,3)p''''\\[9pt] = & q, & q', & q'', & q''', & q'''' \end{array}
\]\endgroup\par% portrait-table end

also
\[
\begin{gathered} q=(a'a'''+b'b''')(a''a''''+b''b'''')a \\[9pt]\quad{}+(a'a'''+b'b''')(a''a''''+b''b'''')bi\\[9pt] \qquad\text{u.\,s.\,f.} \end{gathered}
\]
Es ist dann
\[
\begin{gathered} q'-q=(ba'-ab')(2,4)(b'''-a'''i)\\[9pt]\quad{}=-(ba'-ab')(2,4)p'''i\\[9pt] \qquad\text{u.\,s.\,f.} \end{gathered}
\]



\SourceLeaf{493}

Der Schnitt von $qq'$ mit $oq'''$ hat die complexe Zahl $=$
\[
Q'''=\frac{(1,3)(2,4)(3,0)}{a'''a'''+b'''b'''}p''' \quad\text{u.\,s.\,w.}
\]
Das oben erwähnte Product wird
\[
\begin{gathered}  =-(0,2)(1,3)(2,4)(3,0)(4,1)\\[9pt] Q'''-q=\frac{(1,3)(2,4)(ba'''-ab''')}{a'''a'''+b'''b'''} (b'''-a'''i)\\[9pt] q'-Q'''=\frac{(2,4)(3,0)(a'b'''-b'a''')}{a'''a'''+b'''b'''} (b'''-a'''i) \end{gathered}
\]
\par\medskip\Needspace{5\baselineskip}\pdfbookmark[1]{Art. 3}{pentagramma-art-3}\centerline{[3.]}\nobreak\smallskip
Die Relationen zwischen den Seiten des sphärischen Fünfecks so
\begingroup\small\setlength{\arraycolsep}{3pt}\[
\begin{array}{cccc} \multicolumn{3}{c}{\text{Quadrate der}}&\text{Gleichungen}\\[9pt] \text{Tangenten}&\text{Secanten}&\text{Cosecanten}&\\[9pt] \alpha&\gamma\delta&\frac{\gamma\delta}{\alpha}&1+\alpha=\gamma\delta\\[9pt] \beta&\delta\varepsilon&\frac{\delta\varepsilon}{\beta}&1+\beta=\delta\varepsilon\\[9pt] \gamma&\varepsilon\alpha&\frac{\varepsilon\alpha}{\gamma}&1+\gamma=\varepsilon\alpha\\[9pt] \delta&\alpha\beta&\frac{\alpha\beta}{\delta}&1+\delta=\alpha\beta\\[9pt] \varepsilon&\beta\gamma&\frac{\beta\gamma}{\varepsilon}&1+\varepsilon=\beta\gamma \end{array}
\]\endgroup
Diese Gleichungen sind nicht unabhängig, es ist nemlich identisch:
\[
\begin{gathered} (1+\gamma)(1+\beta-\delta\varepsilon) -(1+\beta)(1+\gamma-\varepsilon\alpha) \\[9pt]\quad{}=\varepsilon\{(1+\beta)\alpha-(1+\gamma)\delta\}\\[9pt] \\[9pt]\quad{}=\varepsilon\{(\alpha\beta-\delta-1)-(\gamma\delta-\alpha-1)\} \end{gathered}
\]
und auf ähnliche Weise wird die fünfte aus dreien der übrigen abgeleitet.
Aus zweien der Grössen $\alpha,\beta,\gamma,\delta,\varepsilon$ folgen die übrigen
\par\nopagebreak% portrait-table start
\begingroup\small\setlength{\arraycolsep}{2pt}\[
\begin{array}{llll} \beta=\frac{1+\alpha+\gamma}{\alpha\gamma}, &\beta=\frac{1+\delta}{\alpha}, &\gamma=\frac{1+\alpha}{\alpha\beta-1}, &\beta=\frac{1+\varepsilon}{\alpha\varepsilon-1}\\[9pt] \delta=\frac{1+\alpha}{\gamma}, &\gamma=\frac{1+\alpha}{\delta}, &\delta=\alpha\beta-1, &\gamma=\alpha\varepsilon-1\\[9pt] \varepsilon=\frac{1+\gamma}{\alpha}, &\varepsilon=\frac{1+\alpha+\delta}{\alpha\delta}, &\varepsilon=\frac{1+\beta}{\alpha\beta-1}, &\delta=\frac{1+\alpha}{\alpha\varepsilon-1} \end{array}
\]\endgroup\par% portrait-table end

Beispiel


\SourceLeaf{494}

\[
\begin{array}{lccrrr} &\cos^{2}&\sin^{2}&&&\\[9pt] \alpha=9&\frac{1}{10}&\frac{9}{10}&71^{\circ}&33'&56''\\[9pt] \beta=\frac{3}{2}&\frac{3}{5}&\frac{2}{5}&39&13&54\\[9pt] \gamma=2&\frac{1}{3}&\frac{2}{3}&54&44&7\\[9pt] \delta=5&\frac{1}{6}&\frac{5}{6}&65&54&19\\[9pt] \varepsilon=\frac{1}{3}&\frac{3}{4}&\frac{1}{4}&30&0&0 \end{array}
\]
Schöne Gleichung
\[
\begin{gathered}3+\alpha+\beta+\gamma+\delta+\varepsilon =\alpha\beta\gamma\delta\varepsilon \\[9pt]\quad{}=\surd(1+\alpha)(1+\beta)(1+\gamma)(1+\delta)(1+\varepsilon)\end{gathered}
\]
\par Der Inhalt des sphärischen Pentagons ist $360^{\circ}$ weniger Summe der Seiten.
Setzt man die Summe $=S$ und
\[
(1+i\surd\alpha)(1+i\surd\beta)(1+i\surd\gamma) (1+i\surd\delta)(1+i\surd\varepsilon)=A+Bi
\]
so wird:
\[
\begin{gathered} A=\alpha\beta\gamma\delta\varepsilon.\cos S\\[9pt] B=\alpha\beta\gamma\delta\varepsilon.\sin S \end{gathered}
\]
\par\medskip\Needspace{5\baselineskip}\pdfbookmark[1]{Art. 4}{pentagramma-art-4}\centerline{[4.]}\nobreak\smallskip
\[
\begin{gathered} G(\alpha x+\beta y+\gamma z)^2 +G'(\alpha'x+\beta'y+\gamma'z)^2 \\[9pt]\quad{}+G''(\alpha''x+\beta''y+\gamma''z)^2\\[9pt] \quad=Axx+Byy+Czz+2ayz+2bxz+2cxy \end{gathered}
\]
\[
\begin{array}{rrrr} (A-G)\alpha+&c\beta+&b\gamma&=0\\[9pt] c\alpha+&(B-G)\beta+&a\gamma&=0\\[9pt] b\alpha+&a\beta+&(C-G)\gamma&=0 \end{array}
\]
\[
\begin{gathered}\{a(A-G)-bc\}\alpha\\[9pt]\quad{}=\{b(B-G)-ac\}\beta\\[9pt]\quad{}=\{c(C-G)-ab\}\gamma\end{gathered}
\]
\[
\frac{bc}{a(A-G)-bc}+\frac{ac}{b(B-G)-ac} +\frac{ab}{c(C-G)-ab}+1=0
\]
\[
\begin{gathered}(A-G)(B-G)(C-G)+2abc\\[9pt]\quad{}=aa(A-G)+bb(B-G)+cc(C-G)\end{gathered}
\]
\[
\alpha\alpha=\frac{1}{1+ \left(\frac{a(A-G)-bc}{b(B-G)-ac}\right)^2+ \left(\frac{a(A-G)-bc}{c(C-G)-ab}\right)^2}
\]

\SourceLeaf{495}

Setzt man
\[
\frac{bc}{a(A-x)-bc}+\frac{ac}{b(B-x)-ac}+\frac{ab}{c(C-x)-ab}+1=\Omega
\]
so wird indefinite
\[
\begin{gathered} \{a(A-x)-bc\}\,\{b(B-x)-ac\}\,\{c(C-x)-ab\}\,\Omega\\[9pt] =-abc(x-G)(x-G')(x-G'') \end{gathered}
\]

Differentiirt man und setzt nach der Differentiation $x=G$, so wird
\[
\begin{gathered} \{a(A-G)-bc\}\,\{b(B-G)-ac\}\{c(C-G)-ab\}\\[9pt] \times abc\Bigl\{\frac{1}{(a(A-G)-bc)^2}+\frac{1}{(b(B-G)-ac)^2}\\[9pt]\quad{}+\frac{1}{(c(C-G)-ab)^2}\Bigr\}\\[9pt] =abc(G-G')(G-G'') \end{gathered}
\]
woraus
\[
\alpha=\surd{\frac{[b(B-G)-ac][c(C-G)-ab]}{[a(A-G)-bc](G-G')(G-G'')}}
\]

Die Schlüsse bedürfen einer Abänderung, wenn eine der Grössen $a,b,c$
verschwindet. Die obige erste Gleichung für $G$ wäre dann eine identische.

\par\medskip\Needspace{5\baselineskip}\pdfbookmark[1]{Art. 5}{pentagramma-art-5}\centerline{[5.]}\nobreak\smallskip
Gleichung der Punkte der Kegelfläche, in welcher die Punkte $(1)(2)(3)(4)(5)$
liegen, Spitze des Kegels im Mittelpunkt der Kugel zugleich Anfangspunkt der
Coordinaten. Achse der $x$ geht durch den Punkt $(3)$, also Ebene der $yz$ geht
durch $(1)$ und $(5)$, Achse der $y$ geht durch $(1)$
\[
\begin{array}{c@{\qquad}c@{\qquad}c@{\qquad}c} &x&y&z\\[9pt] (3)&1&0&0\\[9pt] (4)&\cos1&0&\sin1\\[9pt] (5)&0&\cos3&\sin3\\[9pt] (1)&0&1&0\\[9pt] (2)&\cos5&\cos4&-\cos3\mathord{.}\sin5 \end{array}
\]
Gleichung
\[
(z\mathord{.}\cos1-x\mathord{.}\sin1)(z\mathord{.}\cos3-y\mathord{.}\sin3)\cos2=xy
\]


\SourceLeaf{496}

oder
\[
zz=xz\surd{\alpha}+yz\surd{\gamma}+\frac{1+\alpha+\gamma}{\surd{\alpha\gamma}}xy
\]
Durch Veränderung der Coordinatenflächen lässt sich dieselbe in die Form
bringen
\[
z'z'=Lx'x'+My'y'
\]
Löst man die Gleichung auf
\[
t(2t-1)^2=\alpha\beta\gamma\delta\varepsilon(t-1)
\]
welche eine negative $(G)$ und zwei positive Wurzeln $(G',G'')$ hat, so wird
\[
\begin{gathered} Gz'z'+G'x'x'+G''y'y'=0\\[9pt] GG'G''=-\tfrac14\alpha\beta\gamma\delta\varepsilon\\[9pt] \qquad(G-1)(G'-1)(G''-1)=-\tfrac14\\[9pt] (2G-1)(2G'-1)(2G''-1)=-\alpha\beta\gamma\delta\varepsilon \end{gathered}
\]
Für obiges Beispiel
\[
t(2t-1)^2=20(t-1)
\]
Wurzeln\quad
$-2{,}1973145,\quad+1{,}06931815,\quad+2{,}1279965$

Setzt man
\[
\alpha\beta\gamma\delta\varepsilon=\omega\quad\text{und}\quad \surd{\frac{(18\omega+1)^2}{(3\omega+1)^{3}}}=\cos3\psi
\]
so wird
\[
t=\tfrac13-\tfrac13\cos\psi\mathord{.}\surd{(3\omega+1)}
\]

Das Verhalten der cubischen Gleichung
\[
\frac{t(2t-1)^2}{t-1}=\alpha\beta\gamma\delta\varepsilon
\]
(welche man am bequemsten mit \textsc{Weidenbach}'s Tafel auflöst, wo für
$\frac{1-x}{1+x}=y$
gesucht werden muss $\frac{1}{yxx}=\alpha\beta\gamma\delta\varepsilon$, wonach dann
$2t-1=\frac1x$ wird) übersieht man durch folgende Tafel


\SourceLeaf{497}

\par\nopagebreak% portrait-table start
\begingroup\small\setlength{\arraycolsep}{2pt}\[
\begin{array}{r@{\quad}r@{\quad}r@{\quad}r@{\quad}r@{\quad}r} t&\alpha\beta\gamma\delta\varepsilon&t&\alpha\beta\gamma\delta\varepsilon&t&\alpha\beta\gamma\delta\varepsilon\\[9pt] \infty&\infty&+1.9&+16.6&+1.0&\infty\\[9pt] +10&+400.9&+1.8&+15.2&\mathord{.}\ \mathord{.}\ \mathord{.}\ \mathord{.}\ \mathord{.}&\text{negativ}\\[9pt] +9&+325.1&+1.7&+14.0&-1.0&\infty\\[9pt] +8&+257.1&+1.6&+12.9&-1.618034&+11.0901699\\[9pt] +7&+197.2&+1.5&+12.0&-2&+16.7\\[9pt] +6&+145.2&+1.4&+11.34&-3&+36.7\\[9pt] +5&+101.2&+1.309017&+11.0901699&-4&+64.8\\[9pt] +4&+65.3&+1.3&+11.09&-5&+100.8\\[9pt] +3&+37.5&+1.2&+11.76&-\infty&\infty\\[9pt] +2&+18&+1.1&+15.84&& \end{array}
\]\endgroup\par% portrait-table end


Damit also drei reelle Wurzeln Statt finden, muss
$\alpha\beta\gamma\delta\varepsilon\EqOverGt11.0901699$
oder $\frac{11}{2}+\frac{\surd{125}}{2}$ sein, die Grenzwerthe für $t$ sind also:
$+1.309017=\frac{3+\surd5}{4}$ und $-1.618034=-\frac{1+\surd5}{2}$.

\par\medskip\Needspace{5\baselineskip}\pdfbookmark[1]{Art. 6}{pentagramma-art-6}\centerline{[6].}\nobreak\smallskip
\noindent\hspace*{2em}$A,\,B,\,C,\,D,\,E$\quad die Winkelpunkte des Polygons\par
\noindent\hspace*{2em}$a,\quad b,\quad c,\quad d,\quad e$\quad Pole der Diagonalen\par
\noindent\hspace*{2em}$0,\quad1,\quad2$\quad die drei Hauptachsen entsprechend den Wurzeln
$G,G',G''$ der Gleichung
\[
\begin{gathered}\frac{t(2t-1)^2}{t-1}\\[9pt]\quad{}=(\operatorname{tang}AB\mathord{.}\operatorname{tang}BC\mathord{.}\operatorname{tang}CD\mathord{.}\operatorname{tang}DE\mathord{.}\operatorname{tang}EA)^2\end{gathered}
\]
wo für $G$ die negative Wurzel genommen werden mag; so dass
\[
Guu++G'u'u'+G''u''u''=0
\]
wenn $u$ die Coordinaten irgend eines der Punkte $A,B,C,D,E$ bedeuten,
also zugleich $uu+u'u'+u''u''=1$.

Die Quelle der Hauptsätze ist in den zwei Gleichungen enthalten
\[
\begin{gathered} \text{I.}\qquad\cos0A\mathord{.}\cos0b=-\frac{\operatorname{tang}EA\mathord{.}(2G-1-\operatorname{tang}AB^2)}{4(G-G')(G-G'')}\\[9pt] \text{II.}\qquad\cos0A\mathord{.}\cos0C=-\frac{2G-1-\frac{1}{\sin DE^2}}{4\cos BC\mathord{.}\cos AB(G-G')(G-G'')} \end{gathered}
\]

\SourceLeaf{498}

Die Gleichung I. repräsentirt 30 Gleichungen, die II. hingegen 15 da alle
Permutationen der Achsen und der Winkelpunkte erlaubt sind. Noch zierlicher
(in den anfänglichen Bezeichnungen, wo $\alpha=\operatorname{tang}CD^2$ u.\,s.\,w.)
\begingroup\small\setlength{\arraycolsep}{3pt}\[
\begin{array}{ll} \mathrm{I.}&\left\{\begin{aligned} \cos 0C\mathbin{.}\cos 0d&=\frac{1+\alpha-2G}{4(G'-G)(G''-G)}\mathbin{.}\surd\varepsilon=\mathfrak{A}\mathbin{.}\surd\varepsilon\\[9pt] \cos 0D\mathbin{.}\cos 0c&=\frac{1+\alpha-2G}{4(G'-G)(G''-G)}\mathbin{.}\surd\beta=\mathfrak{A}\mathbin{.}\surd\beta \end{aligned}\right.\\[18pt] \mathrm{II.}&\displaystyle \cos 0B\mathbin{.}\cos 0E=\frac{2\alpha+1-2\alpha G}{4(G'-G)(G''-G)}\mathbin{.}\surd\beta\varepsilon=\mathfrak{a}\mathbin{.}\surd\beta\varepsilon \end{array}
\]\endgroup
Die zehn Gleichungen I. können nur für neun gelten, weil die Multiplication von fünfen dasselbe Resultat gibt wie die Multiplication der fünf übrigen.
Es muss also zwischen den 10 Grössen $\mathfrak{A},\mathfrak{a},\mathfrak{B},\mathfrak{b}$ etc. vier Bedingungsgleichungen geben, welche am zierlichsten so dargestellt werden
\[
\begin{gathered} \beta\mathfrak{b}\mathfrak{c}\mathfrak{C}=\varepsilon\mathfrak{e}\mathfrak{d}\mathfrak{D},\\[9pt] \gamma\mathfrak{c}\mathfrak{d}\mathfrak{D}=\alpha\mathfrak{a}\mathfrak{e}\mathfrak{E}\quad\text{u.\,s.\,w. oder auch}\\[9pt] \beta\mathfrak{a}\mathfrak{A}\mathfrak{C}=\gamma\mathfrak{d}\mathfrak{B}\mathfrak{D},\\[9pt] \gamma\mathfrak{b}\mathfrak{B}\mathfrak{D}=\delta\mathfrak{e}\mathfrak{C}\mathfrak{E}\quad\text{u.\,s.\,w.} \end{gathered}
\]
Es ist aber
\[
\begin{gathered}\cos 0A^2=\frac{\mathfrak{e}\mathfrak{a}\mathfrak{b}}{\mathfrak{c}\mathfrak{d}}\mathbin{.}\alpha =\frac{\mathfrak{D}\mathfrak{b}\gamma}{\mathfrak{E}} =\frac{\mathfrak{C}\mathfrak{c}\delta}{\mathfrak{B}}\quad\text{u.\,s.\,w.,} \\[9pt]\quad{}\qquad\cos 0a^2=\frac{\mathfrak{C}\mathfrak{D}}{\mathfrak{a}}\quad\text{u.\,s.\,w.}\end{gathered}
\]

\par\medskip\Needspace{5\baselineskip}\pdfbookmark[1]{Art. 7}{pentagramma-art-7}\centerline{[7.]}\nobreak\smallskip
1843. April 20. Die excentrischen Anomalien $\varphi,\varphi',\varphi'',\varphi''',\varphi''''$ der Punkte
$A,B,C,D,E$ sind durch die Gleichungen verbunden ($G$ als negativ betrachtet)
\[
\begin{gathered} \frac{\sin\frac12(\varphi'''+\varphi'')}{\cos\frac12(\varphi'''-\varphi'')} =\frac{G}{G''}\mathbin{.}\sin\varphi,\\[9pt]\quad{}\qquad \frac{\cos\frac12(\varphi'''+\varphi'')}{\operatorname{cso}\frac12(\varphi'''-\varphi'')} =\frac{G}{G''}\mathbin{.}\cos\varphi\\[9pt] \frac{\sin\frac12(\varphi'+\varphi'''')}{\cos\frac12(\varphi'-\varphi'''')} =\surd\frac{G(G-1)}{G''(G''-1)}\mathbin{.}\sin\varphi \\[9pt]\quad{}=\frac{G(2G-1}{G''(2G''-1)}\sin\varphi\\[9pt] \frac{\cos\frac12(\varphi'+\varphi'''')}{\cos\frac12(\varphi'-\varphi'''')} =\surd\frac{G(G-1)}{G'(G'-1)}\mathbin{.}\cos\varphi \\[9pt]\quad{}=\frac{G(2G-1)}{G'(2G'-1)}\cos\varphi \end{gathered}
\]
Die Relationen zwischen den Winkeln $\varphi^0,\varphi',\varphi''$ sind am einfachsten auf
folgende Art darzustellen, $\surd\frac{G'}{G'-1}=\xi$, $\surd\frac{G''}{G''-1}=\eta$ gesetzt, wird $\xi\eta=$
\[
\begin{gathered}\frac{\operatorname{tang}\frac12(\varphi'-\varphi'''')}{\operatorname{tang}\frac12(\varphi'''-\varphi'')} \\[9pt]\quad{}=\frac{\operatorname{tang}\frac12(\varphi''-\varphi^0)}{\operatorname{tang}\frac12(\varphi''''-\varphi''')} \\[9pt]\quad{}=\frac{\operatorname{tang}\frac12(\varphi'''-\varphi')}{\operatorname{tang}\frac12(\varphi^0-\varphi'''')} \\[9pt]\quad{}=\frac{\operatorname{tang}\frac12(\varphi''''-\varphi'')}{\operatorname{tang}\frac12(\varphi'-\varphi^0)} \\[9pt]\quad{}=\frac{\operatorname{tang}\frac12(\varphi^0-\varphi''')}{\operatorname{tang}\frac12(\varphi''-\varphi')}\end{gathered}
\]
für $\frac{\xi\eta+1}{\xi\eta-1}$ gibt es einen ähnlichen Ausdruck, der sich hieraus leicht ableiten lässt.


\SourceLeaf{499}

In Zahlen
\par\nopagebreak% portrait-table start
\begingroup\small\relax\setlength{\arraycolsep}{2pt}\[
\begin{array}{r@{\;}c@{\;}rrr@{\quad}rrr@{\quad}rrr@{\quad}rrr} &&&&&&&&&&&\log\operatorname{tang}&\log\operatorname{tang}&\log\Delta\\[9pt] \varphi^0&=&50^0&29'&20''&80^0&54'&55''&49^0&13'&4''&0.79616&0.06418&9.81007\\[9pt] \varphi'&=&92&56&38&55&49&27&15&16&12.5&0.16814&9.43617&9.182\\[9pt] \varphi''&=&162&8&14&83&48&52&59&41&16.5&0.96505&0.23313&9.97901\\[9pt] \varphi'''&=&260&34&22&64&29&16.5&21&13&39&0.32127&9.58931&9.33515\\[9pt] \varphi''''&=&291&6&47&74&57&29&34&35&48&0.57068&9.83871&9. \end{array}
\]\endgroup\par% portrait-table end

\hfill $0.73196=\log\xi\eta$\hspace*{.12\textwidth}\par

\par\medskip\Needspace{5\baselineskip}\pdfbookmark[1]{Art. 8}{pentagramma-art-8}\centerline{[8.]}\nobreak\smallskip
Die $\varphi,\varphi',\varphi'',\varphi''',\varphi''''$ sind nichts anders als die Amplituden zu fünf transscendenten Argumenten, welche um $\frac45K$ zunehmen (in der Bedeutung von \textsc{Jacobi}
p.\,31) und wo der Modulus $k$
\[
\begin{gathered}=\surd\frac{\frac{1}{G'G'}-\frac{1}{G''G''}}{\frac{1}{G'G'}-\frac{1}{GG}} =\sin\mu,\\[9pt]\quad{}\qquad \cos\mu=\surd\frac{\frac{1}{G''G''}-\frac{1}{GG}}{\frac{1}{G'G'}-\frac{1}{GG}}\end{gathered}
\]
Das transcendente Argument selbst, unbestimmt genommen, ist
\[
=\int\frac{x\,\mathrm{d}y-y\,\mathrm{d}x}{\surd(xx+yy)}\mathbin{.} \surd\frac{\frac{1}{G'G'}-\frac{1}{GG}} {\left(\frac{1}{G'G'}-\frac{1}{GG}\right)xx+\left(\frac{1}{G''G''}-\frac{1}{GG}\right)yy}
\]
$\Delta$ in der Bezeichnung von \textsc{Jacobi} gebraucht so dass $\Delta\varphi=\surd(1-kk\sin\varphi^2)$,
es sind
\par\nopagebreak% portrait-table start
\begingroup\small\setlength{\arraycolsep}{2pt}\[
\begin{array}{c|c|c|c|c} & & &\text{propositional}&\\[9pt] \text{die drei Grössen}&\text{ebenso}&&\text{den Zahlen}&\text{oder}\\[9pt] \operatorname{tang}\frac12(\varphi'-\varphi'''')& \operatorname{tang}\frac12(\varphi''-\varphi^0)&\text{u.\,s.\,w.}& G(2G-1)\surd\left(\frac{1}{G'G'}-\frac{1}{GG}\right)& \operatorname{tang}\operatorname{am}\frac45K\\[9pt] \operatorname{tang}\frac12(\varphi'''-\varphi'')& \operatorname{tang}\frac12(\varphi''''-\varphi''')&& -G\surd\left(\frac{1}{G'G'}-\frac{1}{GG}\right)& \operatorname{tang}\operatorname{am}\frac25K\\[9pt] \Delta\varphi^0&\Delta\varphi'&&1&1 \end{array}
\]\endgroup\par% portrait-table end

\par\nopagebreak\begin{center}\rule{.25\textwidth}{.4pt}\end{center}


\clearpage\fancyhead[L]{\small BEMERKUNGEN.}\pdfbookmark[0]{Bemerkungen}{final-bemerkungen}\SourceLeaf{500}
\begin{center}\Large\textbf{BEMERKUNGEN.}\end{center}
\begin{center}\rule{.09\textwidth}{.4pt}\end{center}

In diesem dritten Bande von \textsc{Gauss} Werken habe ich alle Abhandlungen und Aufsätze aus dem
Gebiete der allgemeinen Analysis und zwar speciell aus der Theorie der algebraischen Functionen, der
\textsc{Gauss}sischen Reihen und der Elliptischen Functionen, so wie einige Mittheilungen das \textsc{Pfaff}sche Theorem
und die Construction von Logarithmentafeln betreffend vereinigt. Sie bestehen aus einer Doctordissertation
(früher in Quart gedruckt); aus früher veröffentlichten Abhandlungen: fünf in den \textit{‘Commentationes societatis regiae scientiarum Gottingensis’} (Quart), einer in den \textit{‘Abhandlungen der Königlichen Gesellschaft der}
\textit{Wissenschaften zu Göttingen’} (Quart) und einer im \textsc{Crelle}schen \textit{‘Journal für reine und angewandte Mathematik’} (Quart); ferner aus sechs Anzeigen eigner Abhandlungen in den \textit{‘Göttingischen Gelehrten Anzeigen}
\textit{und Nachrichten’} (Octav); aus zwölf Mittheilungen über nicht eigne Schriften in den \textit{‘Göttingischen gelehrten Anzeigen’} (Octav) (von \textsc{Gauss} nicht unterzeichnet aber in Betreff seiner Autorschaft durch die Acten
der Göttinger Universitäts-Bibliothek verificirt), in der \textit{‘Monatlichen Correspondenz zur Beförderung der Erd-}
\textit{und Himmels-Kunde herausgegeben vom Freiherrn \textsc{von Zach}’} (Octav), in den \textit{‘Astronomischen Nachrichten}
\textit{herausgegeben von \textsc{Schumacher}’} (Quart); und auch aus zwei Erläuterungen in der \textit{‘Monatlichen Correspondenz’}
und in \textsc{Vega}’s \textit{‘Sammlung von Hülfstafeln’} (Quart) zu den \textsc{Gauss}ischen Additions- und Subtractions-Logarithmentafeln. Diese Tafeln selbst habe ich hier nicht abdrucken lassen, weil sie sehr verbreitet sind und
das Format dieser Werke zum Gebrauche solcher Tafeln unbequem sein würde.

Bei der Redaction habe ich dieselben Grundsätze befolgt wie in den frühern Bänden. Zur bessern
Übersicht der Gegenstände in einem so umfangreichen Bande sind die Hauptlehrsätze auf gleiche Weise
durch den Druck ausgezeichnet. Zum leichtern Gebrauch sowohl der ältern Ausgaben als der vorliegenden
habe ich bei den Verweisungen auf Schriften, die nicht in diesem Bande selbst sich finden, statt des Ortes der
Veröffentlichungen die eignen Titel so wie statt der Nummer der Seite die der Artikel angegeben. Auf
Seite 20 in Zeile 15 und 16 ist gemäss einer handschriftlichen Bemerkung \textit{‘objectionem secundam et quar\SourceLeaf{501}%
tam’} statt \textit{‘objectionem tertiam et quartam’} und in Zeile 19 ebenso \textit{‘objectionem tertiam’} statt \textit{‘objectionem}
\textit{primam’} gesetzt. Ausserdem unterscheidet sich die vorliegende Ausgabe von den frühern derselben Schriften nur durch die Berichtigung einiger Druckfehler, wobei ich zum Theil die \textsc{Gauss}ischen’ Original-Manuscripte benutzen konnte.

Aus dem Handschriftlichen Nachlasse habe ich aufgenommen: die den Seiten 30 und 112 beigefügten Noten, den zweiten Theil der Abhandlung \textit{‘Disquisitiones generales circa seriem infinitam etc.’}, eine
ausführliche Abhandlung über Interpolation mit einigen in meinen Bemerkungen Seite 328 u.\,f. erläuterten
Zusätzen, und eine Reihe von Abhandlungen, die sich auf die Elliptischen Functionen beziehen.

Die Handschriftlichen Aufzeichnungen habe ich hier wie auch früher bis auf Berichtigung von unerheblichen Schreibfehlern unverändert abdrucken lassen und meine Einschaltungen, die mit Ausnahme der
\textit{‘Fortsetzung der Untersuchungen über das arithmetisch-geometrische Mittel’} nur in kurzen Sätzen bestehen,
durch [Einklammerung] abgesondert. Die geschichtlichen Angaben und erforderlichen Zusätze für den Nachlass über hypergeometrische Reihen und Interpolation habe ich in den jenen Abhandlungen unmittelbar
folgenden Bemerkungen zusammengestellt.

\bigskip
Die Redaction der \textsc{Gauss}ischen Arbeiten über Elliptische Functionen wurde Seitens der Königlichen
Gesellschaft der Wissenschaften zu Göttingen von \textsc{Riemann} gewünscht und ihm zu dem Zwecke die Handschriften übergeben. Leider hat er weder eine schriftliche noch mündliche Mittheilung aus diesen seinen
Studien hinterlassen. Erst nach dem bedauernswerthen allzufrühen Tode \textsc{Riemann}s und nachdem der vorliegende dritte Band bis auf jenen Theil gedruckt war, konnte ich die \textsc{Gauss}ischen Handschriften zu meiner Bearbeitung übernehmen.

\bigskip
\textsc{Gauss} hat von seinen Untersuchungen der Functionen, die wir jetzt die Elliptischen nennen, nur einen Theil veröffentlicht: eine Anwendung dieser Theorie auf die höhere Arithmetik in der \textit{Summatio quarundam serierum singularium} 1808 \textit{September} und eine Anwendung auf die Bestimmung der Säcularstörungen
der Planeten in der \textit{Determinatio attractionis, quam in punctum quodvis positionis datae exerceret planeta, si}
\textit{ejus massa per totam orbitam ratione temporis, quo singulae partes describuntur, uniformiter esset dipertita,} 1818
\textit{Januar.} Den grössern Theil hat er niedergelegt in einigen unvollständigen Entwürfen zu Anfängen verschiedener Abhandlungen und in zahlreichen zwischen andern Arbeiten sehr zerstreuten Aufzeichnungen einzelner
Formeln. Diese im handschriftlichen Nachlasse befindlichen Untersuchungen habe ich hier in einzelnen
Gruppen zusammengestellt jenachdem sie vom Algorithmus des Aritmetisch-Geometrischen Mittels oder einem anderen diesem analogen und mit diesem in Verbindung gesetzten Algorithmus ausgehen oder aber
sich auf die speciellen Lemniscatischen Functionen beziehen oder endlich die Darstellungen der allgemeinen
Functionen durch Producte als wesentliches Hülfsmittel gebrauchen. Die Untersuchung des \textit{Pentagramma}
\textit{mirificum} bildet eine Anwendung der Fünftheilung der ganzen Elliptischen Integrale.

\bigskip
Die Abhandlung mit den beiden Abschnitten \textit{de origine proprietatibusque generalibus numerorum mediorum arithmetico-geometricorum} und \textit{de functionibus transcendentibus quae ex differentiatione mediorum}
